\section{Conclusion}
\label{sec:conclusion}

This paper completes the BTV accountability programme. Papers~1--3
established that verifiable AI accountability is \emph{technically}
feasible --- guaranteeing construction integrity, persistence
durability, and privacy-preserving maintenance. This paper asks whether
it is \emph{economically} rational to deploy.

Our analysis shows that the cost of non-repudiable transparency scales
inversely with decision volume, reaching \$0.005 per decision at
Enterprise scale ($10^6$ decisions/year) and approaching zero
marginally at Hyperscale. The fixed infrastructure cost of
$\sim$\$5{,}000/year is the dominant term below $\sim$50{,}000
decisions/year and becomes negligible above it.

By introducing a \emph{compliance-credit} mechanism grounded in
existing legal instruments (GDPR Art.~42, AI Act Art.~69), we show
that for operators subject to EU~AI~Act enforcement of high-risk AI
systems, the infrastructure investment breaks even at 500{,}000
decisions/year. This threshold is surpassed by virtually all
production-scale credit-scoring, hiring, and healthcare-triage
systems --- the exact deployments targeted by the regulation.
Enforcement data from 2021--2024 (\$2B+ in SEC record-keeping fines,
\euro{}4.5B+ in GDPR actions~\cite{sec-fines-2024,gdpr-fines-2025})
demonstrate that the penalty exposure the model prices is real, not
hypothetical.

However, our analysis also identifies a market failure for operators
below $\sim$250{,}000 decisions/year, where fixed infrastructure costs
create a Valley of Non-Compliance that rational SMEs may not cross
without structural intervention. This is the primary policy implication
of the work: the technology is ready, and for large operators it is
already economically rational; the remaining obstacle is accessible
infrastructure for the long tail of AI deployments.

Together, the four papers in this programme constitute a vertically
integrated accountability stack: structural correctness at construction
time (Paper~1), verifiable durability at delivery time (Paper~2),
privacy-preserving integrity at maintenance time (Paper~3), and
economic rationality at adoption time (Paper~4). The path to
accountable AI is no longer blocked by scientific unknowns. What
remains is institutional will and the shared infrastructure to lower
the fixed cost of transparency for those who cannot yet afford it.
